\chapter{Funktioner}
\label{chap:ch1}
Funktioner er et centralt begreb i matematikken og danner bro mellem aritmetik og mere avanceret matematisk analyse. En funktion kan opfattes som en særlig sammenhæng mellem to mængder, hvor hvert input (fra definitionsmængden) er knyttet til præcis ét output (fra dispositionsmængden). Når vi forstår funktioner og deres egenskaber, kan vi modellere og løse virkelige problemer med større præcision og fleksibilitet.

I dette kapitel introducerer vi funktionsbegrebet, herunder definitioner, notation og vigtige egenskaber. Derefter undersøger vi almindelige funktionstyper grafisk og udvikler de beslægtede begreber logaritmer og løsning af eksponential- og logaritmeligninger. Disse begreber danner grundlag for mere avancerede matematiske emner senere i bogen.

\section{Funktioners egenskaber}
I matematikken knytter en funktion hvert element i én mængde til et bestemt element i en anden mængde (som kan være den samme mængde). Se for eksempel på et hold i Matematik for softwareingeniører, hvor hver studerende får en karakter fra mængden \(\{12, 10, 7, 4, 02\}\). Antag, at karaktererne er som følger:

\begin{figure}[htbp]
    \centering
    \begin{tikzpicture}[>=Stealth] % Stealth is a nice arrow style
        % Left column of names
        \node (A1) at (0,0) {Alice};
        \node (A2) at (0,-1) {Bob};
        \node (A3) at (0,-2) {Carol};
        \node (A4) at (0,-3) {Dave};
        \node (A5) at (0,-4) {Eve};
        
        % Right column of numbers
        \node (B1) at (4,0) {$12$};
        \node (B2) at (4,-1) {$10$};
        \node (B3) at (4,-2) {$7$};
        \node (B4) at (4,-3) {$4$};
        \node (B5) at (4,-4) {$02$};
        
        % Connecting lines with arrows
        \draw[->, thick, black] (A1.east) -- (B1.west);
        \draw[->, thick, black] (A2.east) -- (B3.west);
        \draw[->, thick, black] (A3.east) -- (B2.west);
        \draw[->, thick, black] (A4.east) -- (B1.west);
        \draw[->, thick, black] (A5.east) -- (B5.west);
    \end{tikzpicture}
    \caption{Eksempel på en funktion, der afbilder navne over i tal.}
    \label{fig:names_to_numbers}
\end{figure}


Denne tildeling af karakterer, som er illustreret i \autoref{fig:names_to_numbers}, er et eksempel på en funktion.

Funktioner spiller en central rolle i matematik og datalogi. De definerer diskrete strukturer som følger og strenge, og de bruges til at analysere algoritmers tidskompleksitet. Mange computerprogrammer er designet til at beregne funktionsværdier. Rekursive funktioner, der er defineret ved hjælp af sig selv, har særlig stor betydning i datalogi. Dette afsnit giver et overblik over de grundlæggende funktionsbegreber, som er nødvendige i matematik for softwareingeniører.

\begin{definition}
Lad $A$ og $B$ være ikke-tomme mængder. En funktion $f$ fra $A$ til $B$ er en tildeling af præcis ét element i $B$ til hvert element i $A$. Vi skriver $f(a)=b$, hvis $b$ er det entydige element i $B$, som funktionen $f$ tildeler elementet $a$ i $A$. Hvis $f$ er en funktion fra $A$ til $B$, skriver vi $f: A \rightarrow B$.
\end{definition}

\begin{remark}
    Funktioner kaldes undertiden også afbildninger eller transformationer.
\end{remark}

Funktioner kan angives på forskellige måder. Nogle gange angiver vi tildelingerne eksplicit, som vist i \autoref{fig:names_to_numbers}. Ofte bruges en forskrift som \(f(x) = x + 1\) til at definere en funktion. I andre tilfælde kan funktionen være givet ved et computerprogram.

\begin{definition}
Hvis $f$ er en funktion fra $A$ til $B$, siger vi, at $A$ er \textbf{definitionsmængden} for $f$, og at $B$ er \textbf{dispositionsmængden} for $f$. Hvis $f(a)=b$, siger vi, at $b$ er \textbf{billedet} af $a$, og at $a$ er et \textbf{urbillede} af $b$. \textbf{Værdimængden}, eller billedmængden, for $f$ er mængden af alle billeder af elementer i $A$. Hvis $f$ er en funktion fra $A$ til $B$, siger vi desuden, at $f$ \textbf{afbilder} $A$ ind i $B$.    
\end{definition}

Når vi definerer en funktion, angiver vi dens definitionsmængde, dens dispositionsmængde og afbildningen af elementerne i definitionsmængden over i dispositionsmængden. To funktioner er ens, hvis de har samme definitionsmængde og samme dispositionsmængde og afbilder hvert element i definitionsmængden over i det samme element i dispositionsmængden. 

Det er vigtigt at bemærke, at en ændring af definitionsmængden eller dispositionsmængden giver en anden funktion. Tilsvarende giver en ændring af afbildningen af elementerne også en anden funktion.


\begin{figure}[htbp]
\centering
\begin{tikzpicture}[>=Stealth, node distance=4cm, thick, scale=1]

    % Draw larger circles for sets A and B
    \draw[thick] (0,0) circle [radius=1.5cm] node at (0, 2) {\Large $A$};
    \draw[thick] (6,0) circle [radius=1.5cm] node at (6, 2) {\Large $B$};

    % Place larger nodes for elements a and b=f(a)
    \node[fill, circle, inner sep=2pt, label=below:{$a$}] (a) at (0,0.3) {};
    \node[fill, circle, inner sep=2pt, label=below:{$b=f(a)$}] (b) at (6,0.3) {};

    % Draw the larger straight arrow representing f
    \draw[->, ultra thick, black] (a) -- (b) node[midway, above] {\Large $f$};

    % Draw the curved arrow representing f at the bottom, touching the circles
    \draw[->, ultra thick, black, bend right=30] ([shift={(270:1.5cm)}]0,0) 
    to node[midway, below] {\Large $f$} 
    ([shift={(270:1.5cm)}]6,0);

\end{tikzpicture}
\caption{En funktion \(f\), der afbilder et element \(a\) i mængden \(A\) over i et element \(b=f(a)\) i mængden \(B\).}
\label{fig:function_mapping}
\end{figure}

De følgende eksempler illustrerer forskellige funktioner. I hvert eksempel beskriver vi definitionsmængden, dispositionsmængden, værdimængden og tildelingen af værdier til elementerne i definitionsmængden.

\begin{example} Hvad er definitionsmængden, dispositionsmængden og værdimængden for den funktion, der tildeler de studerende karakterer, og som er beskrevet i første afsnit af indledningen til dette afsnit?

\begin{solution}
    Lad $G$ være den funktion, der tildeler en studerende på vores hold i matematik for softwareingeniører en karakter. Bemærk for eksempel, at $G$(Alice) $=12$. Definitionsmængden for $G$ er mængden \{Alice, Bob, Carol, Dave, Eve $\}$, og dispositionsmængden er mængden $\{12, 10, 7, 4, 02\}$. Værdimængden for $G$ er mængden $\{12, 10, 7, 02\}$, fordi alle karakterer undtagen $4$ er tildelt mindst én studerende.
\end{solution}
    
\end{example}

\begin{example}
    Lad $f$ være den funktion, der til en bitstreng af længde 2 eller mere tildeler strengens sidste to bit. For eksempel er $f(11010)=10$. Så er definitionsmængden for $f$ mængden af alle bitstrenge af længde 2 eller mere, og både dispositionsmængden og værdimængden er mængden $\{00,01,10,11\}$.
\end{example}

\begin{example}
    Lad $f: \mathbb{Z} \rightarrow \mathbb{Z}$ tildele et helt tal dets kvadrat. Så er $f(x)=x^2$, hvor definitionsmængden for $f$ er mængden af alle hele tal, dispositionsmængden for $f$ er mængden af alle hele tal, og værdimængden for $f$ er mængden af alle hele tal, der er kvadrattal, nemlig $\{0,1,4,9, \ldots\}$.
\end{example}

\subsection*{Injektive og surjektive funktioner}
I matematikken er funktioner et grundlæggende begreb, der bruges til at beskrive sammenhængen mellem to mængder. Men ikke alle funktioner opfører sig ens. For at forstå forskellene introducerer vi begreberne injektive og surjektive funktioner.

Nogle funktioner tildeler aldrig den samme værdi til to forskellige elementer i definitionsmængden. Sådanne funktioner kaldes \textbf{injektive}.

\begin{definition} {Injektive funktioner}
En funktion \( f: A \rightarrow B \) kaldes \textbf{injektiv}, hvis forskellige elementer i \( A \) afbildes over i forskellige elementer i \( B \). Med andre ord: Hvis \( f(a_1) = f(a_2) \), så er \( a_1 = a_2 \). Denne egenskab sikrer, at to forskellige elementer i \( A \) aldrig afbildes over i det samme element i \( B \).     
\end{definition}

For en reel funktion, hvis graf er tegnet i et koordinatsystem, gælder, at funktionen er injektiv, hvis og kun hvis ingen vandret linje skærer grafen mere end én gang. Dette kriterium kaldes den vandrette linjetest.

\begin{definition}{Surjektive funktioner}
En funktion \( f: A \rightarrow B \) kaldes \textbf{surjektiv}, hvis hvert element i \( B \) er billede af mindst ét element i \( A \). Med andre ord findes der for hvert \( b \in B \) mindst ét \( a \in A \), så \( f(a) = b \). Denne egenskab sikrer, at funktionen »dækker« hele mængden \( B \).
\end{definition}

\subsection*{Inverse funktioner}

Se nu på en funktion \( f: A \rightarrow B \), som er både injektiv og surjektiv. Fordi \( f \) er surjektiv, er hvert element i \( B \) billede af et element i \( A \). Fordi \( f \) er injektiv, er hvert element i \( B \) desuden billede af ét entydigt bestemt element i \( A \). Denne entydige korrespondance gør det muligt at definere en ny funktion fra \( B \) til \( A \), som »vender« den afbildning om, der er givet ved \( f \).

\begin{figure}[htbp]
\centering
\begin{tikzpicture}[>=Stealth, thick, scale = 1.5]

% Draw circles representing sets A and B
\draw (0,0) circle (1.5cm);
\draw (6,0) circle (1.5cm);

% Labels for sets A and B
\node at (0,-1) {A};
\node at (6,-1) {B};


% Points in sets with adjusted labels
\node[fill=black, circle, inner sep=1.5pt, label={[xshift=-0.6cm]below:{$a = f^{-1}(b)$}}] (A1) at (0,0.7) {};
\node[fill=black, circle, inner sep=1.5pt, label={[xshift=0.5cm]below:{$b = f(a)$}}] (B1) at (6,0.7) {};


% Arrows for f and f^{-1}
\draw[->] (A1) to[bend left=30] node[midway, above] {$f(a)$} (B1);
\draw[->] (B1) to[bend left=30] node[midway, below] {$f^{-1}(b)$} (A1);

% Arrows for f and f^{-1} below the main circles
\draw[->] (0.6,-1.5) to[bend left=20] node[midway, above] {$f$} (5.4,-1.5);
\draw[->] (5.4,-1.8) to[bend left=20] node[midway, below] {$f^{-1}$} (0.6,-1.8);

\end{tikzpicture}
\caption{Funktionen $f^{-1}$ er den inverse funktion til $f$.}
\label{fig:inv}
\end{figure}

Denne nye funktion kaldes den \textbf{inverse funktion} (eller omvendte funktion) til \( f \) og betegnes \( f^{-1}: B \rightarrow A \). Den inverse funktion \( f^{-1} \) har følgende egenskaber:

\[
f(f^{-1}(b)) = b \quad \text{for alle } b \in B.
\]
\[
f^{-1}(f(a)) = a \quad \text{for alle } a \in A.
\]

Vi kan sammenfatte disse overvejelser i følgende definition

\begin{definition}{Inverse funktioner}
Lad $f$ være en bijektion fra mængden $A$ til mængden $B$. Den inverse funktion til $f$ er den funktion, der til et element $b$ i $B$ tildeler det entydige element $a$ i $A$, for hvilket $f(a)=b$. Den inverse funktion til $f$ betegnes $f^{-1}$. Altså er $f^{-1}(b)=a$, når $f(a)=b$.    
\end{definition}

Disse egenskaber viser, at \( f^{-1} \) reelt ophæver virkningen af \( f \) ved at afbilde hvert element i \( B \) tilbage i det tilsvarende element i \( A \).

\begin{example}


Lad \( f: \mathbb{R} \rightarrow \mathbb{R} \) være defineret ved \( f(x) = 2x + 3 \). Vi kan kontrollere, at \( f \) er både injektiv og surjektiv:

\begin{itemize}
    \item \textbf{Injektiv}: Hvis \( f(x_1) = f(x_2) \), så er \( 2x_1 + 3 = 2x_2 + 3 \). Når vi trækker 3 fra på begge sider, får vi \( 2x_1 = 2x_2 \), og division med 2 giver \( x_1 = x_2 \). Altså er \( f \) injektiv.
    \item \textbf{Surjektiv}: For et vilkårligt \( y \in \mathbb{R} \) kan vi løse \( y = 2x + 3 \) med hensyn til \( x \) og får \( x = \frac{y - 3}{2} \). Da dette \( x \) findes for alle \( y \), er \( f \) surjektiv.
\end{itemize}

Da \( f \) er både injektiv og surjektiv, har den en invers funktion \( f^{-1} \), der er defineret ved

\[
f^{-1}(y) = \frac{y - 3}{2}.
\]

  
\end{example}

\subsection*{Sammensatte funktioner}
I matematikken kan funktioner kombineres til nye funktioner. En vigtig måde at kombinere funktioner på er sammensætning af funktioner. At sammensætte to funktioner består i bund og grund i at anvende den ene funktion på resultaterne af den anden.

\begin{figure}[htbp]
\centering
\begin{tikzpicture}[>=Stealth, thick]

% Draw circles representing sets A, B, and C
\draw (0,0) circle (1.5cm);  % Set A
\draw (5,0) circle (1.5cm);  % Set B
\draw (10,0) circle (1.5cm); % Set C

% Labels for sets A, B, and C
\node at (0,-1) {A};
\node at (5,-1) {B};
\node at (10,-1) {C};

% Points in sets with adjusted labels
\node[fill=black, circle, inner sep=1.5pt, label={[xshift=-0.1cm]below:{$a$}}] (A) at (0,0.7) {};

\node[fill=black, circle, inner sep=1.5pt, label=below:{$g(a)$}] (B) at (5,0.7) {};
\node[fill=black, circle, inner sep=1.5pt, label={[xshift=0.4cm]below:{$f(g(a))$}}] (C) at (10,0.7) {};

% Arrows for g and f
\draw[->] (A) to[bend left=15] node[midway, above] {$g(a)$} (B);
\draw[->] (B) to[bend left=15] node[midway, above] {$f(g(a))$} (C);
\draw[->] ([shift={(360:1.5cm)}]0,0) to node[midway, below] {$g$} (3.5,0);
\draw[->] ([shift={(360:0cm)}]6.5,0) to node[midway, below] {$f$} (8.5,0);

% Arrows for f ∘ g and (f ∘ g)(a)
\draw[->] (A) to[bend left=40] node[midway, above] {$(f \circ g)(a)$} (C);
% Draw the curved arrow representing f at the bottom, touching the circles
\draw[->, bend right=20] ([shift={(270:1.5cm)}]0,0) to node[midway, below] {$f \circ g$} ([shift={(270:1.5cm)}]10,0);

\end{tikzpicture}
\caption{Sammensætningen af $f$ og $g$.}
\label{fig:composite}
\end{figure}

\begin{definition}
    

Lad \( f: B \rightarrow C \) og \( g: A \rightarrow B \) være to funktioner. Den \textbf{sammensatte funktion} af \( f \) og \( g \), som betegnes \( f \circ g \), er en funktion fra \( A \) til \( C \), der er defineret ved

\[
(f \circ g)(x) = f(g(x)),
\]

for alle \( x \in A \).

\end{definition}

Den sammensatte funktion \( f \circ g \) betyder med andre ord, at du først anvender funktionen \( g \) på inputtet \( x \) og derefter anvender funktionen \( f \) på resultatet \( g(x) \). Sammensætning af funktioner er vist i \autoref{fig:composite}.

\begin{example}
    

Se på funktionerne \( f(x) = 2x + 3 \) og \( g(x) = x^2 \). Den sammensatte funktion \( f \circ g \) er givet ved:

\[
(f \circ g)(x) = f(g(x)) = f(x^2) = 2x^2 + 3.
\]

Her kvadrerer funktionen \( g(x) \) inputtet \( x \), hvorefter funktionen \( f(x) \) ganger resultatet med 2 og lægger 3 til.
\end{example}

Lad os nu bytte om på rækkefølgen i sammensætningen og beregne \( g \circ f \):

\[
(g \circ f)(x) = g(f(x)) = g(2x + 3) = (2x + 3)^2.
\]

Bemærk, at \( f \circ g \) og \( g \circ f \) generelt er forskellige funktioner. Det viser, at sammensætning af funktioner ikke er kommutativ.

\begin{example}
    Lad $g$ være funktionen fra mængden $\{a, b, c\}$ til sig selv, så $g(a)=b, g(b)=c$ og $g(c)=a$. Lad $f$ være funktionen fra mængden $\{a, b, c\}$ til mængden $\{1,2,3\}$, så $f(a)=3, f(b)=2$ og $f(c)=1$. Hvad er sammensætningen af $f$ og $g$, og hvad er sammensætningen af $g$ og $f$?

    \begin{solution}
        Sammensætningen $f \circ g$ er defineret ved $(f \circ g)(a)=f(g(a))=f(b)=2$, $(f \circ g)(b)=f(g(b))=f(c)=1$ og $(f \circ g)(c)=f(g(c))=f(a)=3$.
        
        Bemærk, at $g \circ f$ ikke er defineret, fordi værdimængden for $f$ ikke er en delmængde af definitionsmængden for $g$.
    \end{solution}
\end{example}


\begin{example}
Lad $f$ og $g$ være funktionerne fra mængden af hele tal til mængden af hele tal, der er defineret ved $f(x)=2 x+3$ og $g(x)=3 x+2$. Hvad er sammensætningen af $f$ og $g$? Hvad er sammensætningen af $g$ og $f$?

\begin{solution}
Både sammensætningen $f \circ g$ og sammensætningen $g \circ f$ er defineret. Desuden er
\[
(f \circ g)(x)=f(g(x))=f(3 x+2)=2(3 x+2)+3=6 x+7
\]
og
\[
(g \circ f)(x)=g(f(x))=g(2 x+3)=3(2 x+3)+2=6 x+11 .
\]
\end{solution}
\end{example}

\begin{remark}
Selv om \( f \circ g \) og \( g \circ f \) begge er defineret for funktionerne \( f \) og \( g \) i det foregående eksempel, er de ikke ens. Sammensætning af funktioner er med andre ord generelt ikke kommutativ.
  
\end{remark}

Når en funktion sammensættes med sin inverse funktion, uanset i hvilken rækkefølge, får man en identitetsfunktion. Antag for at se dette, at \( f \) er en bijektion fra mængden \( A \) til mængden \( B \). Så findes den inverse funktion \( f^{-1} \), og den er en bijektion fra \( B \) til \( A \). Den inverse funktion vender den oprindelige funktions korrespondance om, så \( f^{-1}(b)=a \), når \( f(a)=b \), og \( f(a)=b \), når \( f^{-1}(b)=a \).

Derfor er

\[
\left(f^{-1} \circ f\right)(a) = f^{-1}(f(a)) = f^{-1}(b) = a,
\]

og

\[
\left(f \circ f^{-1}\right)(b) = f\left(f^{-1}(b)\right) = f(a) = b.
\]

Følgelig er \( f^{-1} \circ f = I_A \) og \( f \circ f^{-1} = I_B \), hvor \( I_A \) og \( I_B \) er identitetsfunktionerne på henholdsvis mængden \( A \) og mængden \( B \). Det vil sige, at \( \left(f^{-1}\right)^{-1} = f \).

\begin{example}
    Hvis $f: \mathbb{R} \rightarrow \mathbb{R}$ er defineret ved $f(x)=2 x+3$, så er $f^{-1}(x)=\frac{x-3}{2}$. Den sammensatte funktion $f \circ f^{-1}$ er da identitetsfunktionen $I_{\mathbb{R}}$ på de reelle tal, dvs. $f\left(f^{-1}(x)\right)=x$ for alle $x \in \mathbb{R}$.
\end{example}

\section{Grafisk identifikation af funktionstyper}
Det er grundlæggende i matematikken at forstå, hvordan forskellige typer af funktioner opfører sig. Funktioner kan klassificeres ud fra mønstrene i deres grafer, som giver værdifuld indsigt i deres egenskaber. I dette afsnit undersøger vi forskellige typer af funktioner, herunder lineære funktioner, andengradsfunktioner, eksponentialfunktioner med flere. Ved at undersøge deres grafer kan vi identificere vigtige træk som skæringspunkter med akserne, hældninger, krumning og asymptotisk opførsel, så vi effektivt kan skelne mellem de forskellige funktionstyper.

\subsection*{Lineære funktioner}
Den generelle ligning for en lineær funktion er givet ved

\[
y = ax + b \quad \text{(ofte skrevet som } y = mx + b\text{)},
\]

hvor \(a\) (eller \(m\)) er hældningen, og \(b\) er skæringen med y-aksen. Definitionsmængden for funktionen er alle reelle tal. Ligningen står på hældning-skæringsform, fordi \(a\) (eller \(m\)) angiver hældningen, og \(b\) angiver skæringen med y-aksen. Hvis \(a = 0\), reduceres funktionen til \(y = b\), som er en konstant funktion.

Grundfunktionen for en lineær ligning er

\[
y = x.
\]

Den transformerede funktion kan skrives på punkt-hældningsform som

\[
y = y_1 + a(x - x_1),
\]

hvor grafen går gennem punktet \((x_1, y_1)\) og har hældningen \(a\). I denne form gælder:
\begin{itemize}
    \item \(a\) er den lodrette skalering (hældningen),
    \item \(y_1\) angiver den lodrette forskydning,
    \item \(x_1\) angiver den vandrette forskydning.
\end{itemize}

Punkt-hældningsformen kan også skrives som

\[
y - y_1 = a(x - x_1),
\]

hvor koordinaterne til det faste punkt \((x_1, y_1)\) optræder med negativt fortegn. Formen \(y = y_1 + a(x - x_1)\) udtrykker \(y\) eksplicit ved \(x\), hvilket gør den lettere at indtaste på en grafregner.

\begin{figure}[htbp]
    \centering
    \resizebox{\textwidth}{!}{\DaLinearPanels}
    \caption{Lineære funktioner}
    \label{fig:book_image}
\end{figure}

Grafen for en lineær funktion er en ret linje. Grundfunktionen \(y = x\) er vist til venstre i \autoref{fig:book_image}, hældning-skæringsformen i midten og punkt-hældningsformen til højre.

For hældning-skæringsformen gælder: »Start i \(b\) på \(y\)-aksen, gå \(x\) enheder vandret, og stig \(ax\) enheder lodret.« For punkt-hældningsformen gælder: »Start i \((x_1, y_1)\), gå \((x - x_1)\) enheder vandret, og stig \(a(x - x_1)\) enheder lodret.«

\subsection*{Andengradsfunktioner}
Den generelle ligning for en andengradsfunktion er givet ved

\[
y = ax^2 + bx + c,
\]

hvor \(a \neq 0\), og \(a\), \(b\) og \(c\) er konstanter. Definitionsmængden for funktionen er alle reelle tal.

\begin{figure}[htbp]
    \centering
    \resizebox{\textwidth}{!}{\DaQuadraticPanels}
    \caption{Andengradsfunktioner}
    \label{fig:book_image2}
\end{figure}

Grundfunktionen for en andengradsligning er

\[
y = x^2,
\]

hvor parablens toppunkt ligger i origo \((0, 0)\).

Den transformerede funktion kan skrives på toppunktsform som

\[
y = k + a(x-h)^2,
\]

hvor parablens toppunkt ligger i \((h, k)\). I denne form gælder:
\begin{itemize}
    \item \(k\) angiver den lodrette forskydning,
    \item \(h\) angiver den vandrette forskydning,
    \item \(a\) angiver den lodrette skalering.
\end{itemize}

Toppunktsformen kan også skrives som

\[
y - k = a(x-h)^2,
\]

men når \(y\) udtrykkes eksplicit ved \(x\), er ligningen lettere at indtaste på en grafregner.

Grafen for en andengradsfunktion er en parabel (af det græske ord for »langs en bolds bane«). Parablen er konveks (krummer opad), hvis \(a > 0\), og konkav (krummer nedad), hvis \(a < 0\). Denne opførsel er illustreret i \autoref{fig:book_image2}.

\subsection*{Potensfunktioner}
Den generelle ligning for en potensfunktion er givet ved

\[
y = ax^b,
\]

hvor \(a\) og \(b\) er konstanter forskellige fra nul. Funktionens definitionsmængde afhænger af værdien af \(b\):
\begin{itemize}
    \item Hvis \(b > 0\), er definitionsmængden alle reelle tal.
    \item Hvis \(b < 0\), udelukker definitionsmængden \(x = 0\) for at undgå division med nul.
    \item Hvis \(b\) ikke er et helt tal, udelukker definitionsmængden som regel negative tal for at undgå at tage rødder af negative tal.
\end{itemize}
I de fleste anvendelser er definitionsmængden begrænset til ikke-negative tal.

\begin{figure}[htbp]
    \centering
    \resizebox{\textwidth}{!}{\DaPowerPanels}
    \caption{Potensfunktioner}
    \label{fig:book_image3}
\end{figure}

Grundfunktionen for en potensfunktion er

\[
y = x^b.
\]

For den generelle potensfunktion \(y = ax^b\) gælder:
\begin{itemize}
    \item Hvis \(b > 0\), varierer \(y\) ligefrem med \(b\)'te potens af \(x\), dvs. \(y\) er ligefrem proportional med \(b\)'te potens af \(x\).
    \item Hvis \(b < 0\), varierer \(y\) omvendt med \(b\)'te potens af \(x\), dvs. \(y\) er omvendt proportional med \(b\)'te potens af \(x\).
\end{itemize}

Skaleringsfaktoren \(a\) fungerer som proportionalitetskonstant.

Den forskudte form af en potensfunktion er

\[
y = d + a(x - c)^b,
\]

hvor \(c\) og \(d\) er henholdsvis den vandrette og den lodrette forskydning. Dette kan sammenlignes med de forskudte former af lineære funktioner og andengradsfunktioner:

\[
\begin{aligned}
& y = y_1 + a(x - x_1) &\text{(lineær funktion)}, \\
& y = k + a(x - h)^2 &\text{(andengradsfunktion)}.
\end{aligned}
\]

Medmindre andet er angivet, betyder »potensfunktion« den ikke-forskudte form, \(y = ax^b\).

\autoref{fig:book_image3} viser graferne for potensfunktioner for forskellige værdier af \(b\). I alle tilfælde er \(a > 0\). Grafens form og krumning afhænger af værdien af \(b\):
\begin{itemize}
    \item Hvis \(b > 0\), går grafen gennem origo.
    \item Hvis \(b < 0\), har grafen akserne som asymptoter.
    \item Funktionen er voksende, hvis \(b > 0\), og aftagende, hvis \(b < 0\).
    \item Grafen er konveks (krummer opad), hvis \(b > 1\) eller \(b < 0\), og konkav (krummer nedad), hvis \(0 < b < 1\).
\end{itemize}
Grafens krumning beskriver, hvor hurtigt \(y\) vokser. For \(b > 0\) betyder konveks, at \(y\) vokser med en stigende hastighed, mens konkav betyder, at \(y\) vokser med en aftagende hastighed.

\subsection*{Eksponentialfunktioner}

Den generelle ligning for en eksponentialfunktion er givet ved

\[
y = a b^x,
\]

hvor \(a\) og \(b\) er konstanter, \(a \neq 0\), \(b > 0\) og \(b \neq 1\). Definitionsmængden for funktionen er alle reelle tal.

Grundfunktionen for en eksponentialligning er

\[
y = b^x,
\]
hvor asymptoten er \(x\)-aksen.

\begin{figure}[htbp]
    \centering
    \resizebox{\textwidth}{!}{\DaExponentialPanels}
    \caption{Eksponentialfunktioner}
    \label{fig:book_image4}
\end{figure}

I ligningen \(y = a b^x\) siger vi, at »\(y\) varierer eksponentielt med \(x\)«. Det betyder, at \(y\) ændres med en konstant faktor \(b\), hver gang \(x\) øges med én enhed.

Den forskudte form af eksponentialfunktionen er

\[
y = a b^x + c,
\]

hvor asymptoten er linjen \(y = c\). Medmindre andet er angivet, betyder »eksponentialfunktion« den ikke-forskudte form \(y = a b^x\).

\autoref{fig:book_image4} illustrerer eksponentialfunktioner for forskellige værdier af \(a\) og \(b\). Grafens vigtigste egenskaber er:
\begin{itemize}
    \item Konstanten \(a\) er grafens skæring med \(y\)-aksen.
    \item Funktionen er voksende, hvis \(b > 1\), og aftagende, hvis \(0 < b < 1\), forudsat at \(a > 0\).
    \item Hvis \(a < 0\), er funktionens opførsel omvendt: Den er aftagende, hvis \(b > 1\), og voksende, hvis \(0 < b < 1\).
    \item Grafen er konveks (krummer opad), hvis \(a > 0\), og konkav (krummer nedad), hvis \(a < 0\).
\end{itemize}

Matematikere bruger ofte en af to særlige konstanter som grundtal for en eksponentialfunktion: enten 10, som er grundtallet i titalssystemet, eller det naturligt forekommende tal \( e \), som er tilnærmelsesvis 2.71828. Disse grundtal har stor betydning i mange matematiske anvendelser.

\begin{definition}{Særlige eksponentialfunktioner}

\begin{align*}
y &= a \cdot 10^{bx} \hspace{1cm} \text{eksponentialfunktion med grundtal 10} \\
y &= a \cdot e^{bx}  \hspace{1.2cm} \text{naturlig eksponentialfunktion (grundtal }e\text{),}
\end{align*}

hvor $a$ og $b$ er konstanter, og definitionsmængden er alle reelle tal.
    
\end{definition}

For at generalisere eksponentialfunktionen ganges variablen i eksponenten ofte med en konstant. De (ikke-forskudte) generelle former af disse eksponentialfunktioner er givet nedenfor:

\[
y = a \cdot 10^{bx}
\quad \text{og} \quad
y = a \cdot e^{bx}
\]

Funktionerne kan generaliseres yderligere ved at indføre forskydninger i både \(x\)- og \(y\)-retningen. De forskudte former er:

\[
y = a \cdot 10^{b(x - c)} + d
\quad \text{og} \quad
y = a \cdot e^{b(x - c)} + d
\]

Eksponentialfunktionen med grundtal \(e\) har især en betydelig fordel i differentialregning, fordi ændringshastigheden for \(e^x\) er lig med \(e^x\) selv.

\section{Logaritmer}
Ethvert positivt tal kan skrives som en potens af 10. For eksempel er

\[
\begin{aligned}
3 & = 10^{0.477 \ldots} \\
5 & = 10^{0.6989 \ldots} \\
15 & = 10^{1.1760 \ldots}
\end{aligned}
\]

Eksponenterne \(0.4771 \ldots\), \(0.6989 \ldots\) og \(1.1760 \ldots\) kaldes 10-talslogaritmerne af henholdsvis 3, 5 og 15:

\[
\begin{aligned}
\log 3 & = 0.4771 \ldots \\
\log 5 & = 0.6989 \ldots \\
\log 15 & = 1.1760 \ldots
\end{aligned}
\]

For bedre at forstå, hvad logaritmer betyder, kan du trykke \texttt{LOG} 3 på din lommeregner. Du får:

\[
\log 3 = 0.477121254 \ldots
\]

Opløft derefter 10 til denne potens uden at afrunde. Du får:

\[
10^{0.477121254 \ldots} = 3
\]

Potenser af 10 har de sædvanlige egenskaber for potenser. For eksempel er
\[
\begin{aligned}
15 & = (3)(5) = \left(10^{0.4771 \ldots}\right)\left(10^{0.6990 \ldots}\right) \\
& = 10^{0.4771 \ldots + 0.6990 \ldots} \\
& = 10^{1.1761 \ldots}
\end{aligned}
\]

Det betyder, at \(10^{0.4771 \ldots + 0.6990 \ldots} = 10^{1.1761 \ldots}\). Her lægger du eksponenterne sammen og beholder det samme grundtal. Du kan kontrollere med din lommeregner, at \(10^{1.1761 \ldots}\) er tilnærmelsesvis 15; de viste decimaler er afrundet.

Af dette eksempel kan du slutte, at logaritmer har de samme egenskaber som eksponenter. Det er forventeligt, fordi logaritmer \textit{er} eksponenter. For eksempel er

\[
\log(3 \cdot 5) = \log 3 + \log 5 \quad \text{\textit{Logaritmen af et produkt er summen af faktorernes logaritmer}.}
\]

Ud fra de tidligere angivne værdier kan du også vise, at:

\[
\log \frac{15}{3} = \log 15 - \log 3 \quad \text{\textit{Logaritmen af en kvotient}.}
\]

Denne egenskab er rimelig, fordi man dividerer potenser med samme grundtal ved at trække eksponenterne fra hinanden:

\[
\frac{15}{3} = \frac{10^{1.1760 \ldots}}{10^{0.477 \ldots}} = 10^{1.1760 \ldots - 0.4771 \ldots} = 10^{0.6989 \ldots} = 5
\]

Da en potens kan skrives som et produkt, kan du finde logaritmen af en potens på følgende måde:

\[
\begin{aligned}
\log 3^4 & = \log (3 \cdot 3 \cdot 3 \cdot 3) = \log 3 + \log 3 + \log 3 + \log 3 \\
& = 4 \log 3 \quad \text{\textit{Saml ensartede led}.}
\end{aligned}
\]

Logaritmen af en potens er eksponenten gange logaritmen af det tal, der opløftes til potensen. For at kontrollere resultatet kan du bemærke, at \(3^4 = 81\). Tryk \(4 \times \texttt{LOG}\,3\) på din lommeregner, og du vil se, at det giver \(1.9084 \ldots\).

\begin{definition}{10-talslogaritmer}

\[
\log x=y \iff 10^y=x, \qquad x>0
\]

\textit{Med ord}: $\log x$ er eksponenten i den potens af 10, der giver $x$   
\end{definition}

Ordet logaritme kommer af de græske ord \emph{logos}, der betyder »forhold«, og \emph{arithmos}, der betyder »tal«. Før lommeregneren blev opfundet, blev 10-talslogaritmer beregnet tilnærmelsesvis ved hjælp af uendelige rækker og opført i tabeller. Produkter med mange faktorer, som
\[
(357)(4.367)(22.4)(3.142)
\]
kunne beregnes ved at lægge faktorernes logaritmer sammen i stedet for møjsommeligt at gange flere talpar sammen. John Napier introducerede logaritmerne i begyndelsen af det syttende århundrede, og Henry Briggs udviklede og udgav efterfølgende omfattende tabeller over titalslogaritmer (med grundtal 10).

\begin{custombox}{Regneregler for 10-talslogaritmer}
For $x>0$, $y>0$ og $r\in\mathbb{R}$ gælder følgende regler:
\begin{itemize}
    \item Logaritmen af et produkt:
    
    \[
    \log x y=\log x+\log y
    \]
    
    \textit{Med ord}: $\log$ af et produkt er summen af faktorernes log.
    \vspace{0.2cm}
    \item Logaritmen af en kvotient:
    
    \[
    \log \frac{x}{y}=\log x-\log y
    \]
    \vspace{0.1cm}
    \textit{Med ord}: $\log$ af en kvotient er tællerens log minus nævnerens $\log$.
    \vspace{0.2cm}
    \item Logaritmen af en potens:
    \[
    \log x^r=r \log x
    \]
    \textit{Med ord}: $\log$ af en potens er eksponenten gange log af det tal, der opløftes til potensen.
\end{itemize}

   
\end{custombox}

\begin{example} Find $x$, hvis $\log_{10} 10^{3.721}=x$

\begin{solution}
Ifølge definitionen er logaritmen eksponenten til 10. Altså er $x=3.721$.
\end{solution}
\end{example}

\begin{example} Find $x$, hvis $0.258=10^x$

\begin{solution}
   Ifølge definitionen er $x$, eksponenten til 10, logaritmen af 0.258.

    \[
    x=\log_{10} 0.258=-0.5883 \ldots
    \]
\end{solution}
    
\end{example}

\begin{custombox}{Det vigtigste, du skal huske om logaritmer, er dette}
    \textbf{En logaritme er en eksponent.}
\end{custombox}

\subsection*{Logaritmer med vilkårligt grundtal: reglen for basisskift}
Hvis \(x = 10^y\), er \(y\) 10-talslogaritmen af \(x\). Tilsvarende gælder, at hvis \(x = 2^y\), er \(y\) 2-talslogaritmen af \(x\). Den eneste forskel på disse logaritmer er det tal, der fungerer som grundtal. For at skelne mellem logaritmer med forskellige grundtal skrives grundtallet som et sænket tegn efter forkortelsen »log«. For eksempel:

\[
\begin{aligned}
3 &= \log_2 8 \Leftrightarrow 2^3 = 8, \\
4 &= \log_3 81 \Leftrightarrow 3^4 = 81, \\
2 &= \log_{10} 100 \Leftrightarrow 10^2 = 100.
\end{aligned}
\]

Symbolet \(\log_2 8\) udtales »logaritmen med grundtal 2 af 8«. Symbolet \(\log_{10} 100\) er naturligvis det samme som \(\log 100\), som defineret i det foregående afsnit. Bemærk, at en logaritme i alle tilfælde er en eksponent.

\begin{definition}{Logaritme med vilkårligt grundtal}
\textit{Algebraisk}:

\hspace{1cm}$\log _b x=y$ hvis og kun hvis $b^y=x, \quad$ hvor $b>0, b \neq 1$, og $x>0$

\vspace{0.3cm}

\textit{Med ord}:

\hspace{1cm} $\log _b x=y$ betyder, at $y$ er den eksponent til $b$, der giver $x$ som resultat.
    
\end{definition}

Måden, du udtaler logaritmesymbolet på, giver dig en måde at huske definitionen på. De næste to eksempler viser, hvordan du gør.

\begin{example} Skriv $\log _5 c=a$ på eksponentiel form.

\begin{solution}

Tænk sådan:
\begin{itemize}
    \item »\(\log_5 \ldots\)« læses »logaritmen med grundtal 5 \(\ldots\)«, dvs. 5 er grundtallet.
    \item En logaritme er en eksponent. Da \(\log\) er lig med \(a\), må \(a\) være eksponenten.
    \item Det »resultat«, man får af \(5^a\), er logaritmens argument, som betegnes \(c\).
\end{itemize}

Skriv kun dette: 

\[
5^a = c
\]
\end{solution}
\end{example}

\begin{example} Antag, at $z>0$, $z\neq 1$ og $m>0$, og skriv $z^4=m$ på logaritmisk form.

\begin{solution}
    $\log _z m=4$
\end{solution}
    
\end{example}

To grundtal for logaritmer bruges så ofte, at de har deres egen tast på de fleste lommeregnere. Det ene er 10-talslogaritmen, også kaldet titalslogaritmen, som blev behandlet i det foregående afsnit. Det andet er logaritmen med grundtal \(e\), kaldet den naturlige logaritme, hvor \(e = 2.71828 \ldots\) er et naturligt forekommende tal (ligesom \(\pi\)), som vil være en fordel i dine videre matematikstudier.

Symbolet \(\ln x\) (udtales »el-en af \(x\)«) bruges om naturlige logaritmer og er defineret som:
\[
\ln x = \log_e x
\]

\begin{definition}{Titalslogaritme og naturlig logaritme}
\hspace{1cm} \textit{Titalslogaritme}: Symbolet $\log x$ betyder $\log _{10} x$.

\hspace{1cm}  \textit{Naturlig}: \hspace{0.2cm}Symbolet $\ln x$ betyder $\log _e x$, hvor $e$ er en konstant lig med $2.71828182845 \ldots$
\end{definition}

\begin{example}\label{ex:log5-17} Find $\log _5 17$. Kontrollér dit svar med en passende numerisk metode.

\begin{solution} Lad $x=\log _5 17$.

\begin{equation*}
\begin{aligned}
& 5^x=17 \\
& \log _{10} 5^x=\log _{10} 17 \\
& x \log _{10} 5=\log _{10} 17 \\
& x=\frac{\log _{10} 17}{\log _{10} 5}=1.7603 \ldots \\
& \log _5 17=1.7603 \ldots \\
& 5^{1.7603 \ldots}=17
\end{aligned}
\end{equation*}
\end{solution}
\end{example}

Bemærk i dette eksempel, at 5-talslogaritmen af et tal er ligefrem proportional med 10-talslogaritmen af samme tal. Eksemplets konklusion kan udtrykkes således:

\[
\log_5 17 = \frac{1}{\log_{10} 5} \cdot \log_{10} 17 = 1.4306 \ldots \log_{10} 17
\]

For at finde 5-talslogaritmen af et vilkårligt tal skal du blot gange tallets 10-talslogaritme med \(1.4306 \ldots\) (dvs. dividere med \(\log_{10} 5\)).

Denne proportionale sammenhæng kaldes reglen for basisskift. Ud fra resultaterne i \autoref{ex:log5-17} kan du skrive:

\[
\log_5 17 = \frac{\log_{10} 17}{\log_{10} 5}
\]

Bemærk, at logaritmen med det ønskede grundtal står isoleret på ligningens venstre side, mens de to logaritmer på højre side har samme grundtal -- typisk et, som findes på din lommeregner. Boksen nedenfor viser reglen for grundtallene \(a\) og \(b\) med argumentet \(x\):

\begin{custombox}{Reglen for basisskift for logaritmer}

\begin{equation*}
\log _a x=\frac{\log_b x}{\log_b a} \quad \text { eller } \quad \log_a x=\frac{1}{\log_b a}\left(\log_b x\right)
\end{equation*}

hvor $a>0$, $a\neq 1$, $b>0$, $b\neq 1$ og $x>0$.
    
\end{custombox}

\begin{example}
Find $\ln 29$ ved hjælp af reglen for basisskift med 10-talslogaritmer. Kontrollér dit svar direkte ved at trykke $\ln 29$ på din lommeregner.

\begin{solution}

\[
\ln 29=\frac{\log 29}{\log e}=\frac{1.4623 \ldots}{0.4342 \ldots}=3.3672 \ldots
\]

\hspace{0.9cm} Direkte: $\quad \ln 29=3.3672 \ldots,$

hvilket stemmer med det svar, vi fik ved hjælp af reglen for basisskift.
\end{solution}
\end{example}

\begin{custombox}{Logaritmeregler}
\setlength{\leftskip}{1cm}  % Reset the text indent
\setlength{\rightskip}{1cm} % Reset the right indent
Lad $b>0$ med $b\neq 1$, og lad $x>0$, $y>0$ og $r\in\mathbb{R}$.

Logaritmen af en potens:
$$
\log _b x^r=r \log _b x
$$

\textit{Med ord}: Logaritmen af en potens er produktet af eksponenten og logaritmen af det tal, der opløftes til potensen.
\vspace{0.5cm}
Logaritmen af et produkt:
$$
\log _b(x y)=\log _b x+\log _b y
$$
\textit{Med ord}: Logaritmen af et produkt er summen af faktorernes logaritmer.
\vspace{0.5cm}
Logaritmen af en kvotient:
$$
\log _b \frac{x}{y}=\log _b x-\log _b y
$$

\textit{Med ord}: Logaritmen af en kvotient er tællerens logaritme minus nævnerens logaritme.

\setlength{\leftskip}{0cm}  % Reset the text indent
\setlength{\rightskip}{0cm} % Reset the right indent

\end{custombox}

\subsection*{Løsning af eksponential- og logaritmeligninger}
Logaritmer giver en metode til at løse en ligning med en variabel i eksponenten eller en ligning, der allerede indeholder logaritmer. Det viser vi i de næste eksempler.

\begin{example} Løs eksponentialligningen $7^{3 x}=983$ algebraisk ved hjælp af logaritmer.

\begin{solution}
\begin{align*}
& 7^{3x} = 983  \\
& \log 7^{3x} = \log 983 && \texttt{\small Tag 10-talslogaritmen på begge sider.} \\
& 3x \log 7 = \log 983 && \texttt{\small Brug reglen for logaritmen af en potens.} \\
& x = \frac{\log 983}{3 \log 7} && \texttt{\small Dividér begge sider med koefficienten til x.}\\
& x = 1.1803 \ldots   
\end{align*}


\end{solution}
    
\end{example}

\begin{example} Løs ligningen
    \begin{equation*}
\log _2(x-1)+\log _2(x-3)=3
\end{equation*}

\newpage

\begin{solution}

\begin{align*}
& \log _2(x-1)+\log _2(x-3)=3  \\
& \log _2[(x-1)(x-3)]=3 && \texttt{\small Brug reglen for logaritmen af et produkt.} \\
& 2^3=(x-1)(x-3) && \texttt{\small Brug definitionen af logaritmen.} \\
& 8=x^2-4 x+3 && \texttt{\small Gang parenteserne ud.}\\
& x^2-4 x-5=0 && \texttt{\small Få nul på den ene side. Brug lighedens }\\
&\vspace{-0.2cm}   && \texttt{\small symmetriegenskab.} \\
& (x-5)(x+1)=0 && \texttt{\small Løs ved faktorisering.} \\
& x=5 \quad \text{ eller } \quad x=-1
\end{align*}
Vi skal være forsigtige her, fordi løsningerne i sidste trin er løsninger til andengradsligningen, og vi skal sikre os, at de også er løsninger til den oprindelige logaritmeligning. Kontrollér ved at indsætte løsningerne i den oprindelige ligning.

\begin{minipage}[t]{0.45\textwidth}
Hvis $x=5$, så er
\[
\begin{aligned}
& \log _2(5-1)+\log _2(5-3) \\
& =\log _2 4+\log _2 2 \\
& =2+1=3
\end{aligned}
\]
\end{minipage}%
\hfill
\begin{minipage}[t]{0.45\textwidth}
Hvis $x=-1$, så er
\[
\begin{aligned}
& \log _2(-1-1)+\log _2(-1-3) \\
& =\log _2(-2)+\log _2(-4)
\end{aligned}
\]
hvilket ikke er defineret.
\end{minipage}

\end{solution}

\end{example}

\begin{example} Løs ligningen

\[
e^{2 x}-3 e^x+2=0
\]

\begin{solution}

\[
\begin{aligned}
& e^{2 x}-3 e^x+2=0 \\
& \left(e^x\right)^2-3 e^x+2=0
\end{aligned}
\]
Vi indser, at dette er en andengradsligning i variablen $e^x$. Med løsningsformlen for andengradsligningen får du

\[
\begin{aligned}
& e^x=\frac{+3 \pm \sqrt{9-4(2)}}{2}=\frac{3 \pm 1}{2} \\
& e^x=2 \text { eller } e^x=1
\end{aligned}
\]


Nu skal du løse disse to ligninger.
\[
\begin{array}{ll}
e^x=2 & \qquad e^x=1 \\
x=\ln 2=0.6931 \ldots & \qquad x=0
\end{array}
\]
Kontrol:

\[
\begin{array}{ll}
e^{2 \ln 2}-3 e^{\ln 2}+2 & \qquad \left(e^0\right)^2-3 e^0+2 \\
=\left(e^{\ln 2}\right)^2-3 e^{\ln 2}+2 & \qquad =1^2-3(1)+2=0 \\
=2^2-3(2)+2=0 &
\end{array}
\]
Begge løsninger er korrekte.
\end{solution}  
\end{example}

\newpage

\begin{example} Løs logaritmeligningen $\ln (x+3)+\ln (x+5)=0$

\begin{solution}
    \[
\begin{aligned}
& \ln (x+3)+\ln (x+5)=0 \\
& \ln [(x+3)(x+5)]=0 \\
& (x+3)(x+5)=e^0=1 \\
& x^2+8 x+15=1 \\
& x^2+8 x+14=0 \\
& x=-2.5857 \ldots \quad \text { eller } \qquad x=-5.4142 \ldots
\end{aligned}
\]
Kontrol:
\[
\begin{aligned}
& x=-2.5857 \ldots: \\
& \ln (-2.5857 \ldots+3)+\ln (-2.5857 \ldots+5) \\
& =\ln (0.4142 \ldots)+\ln (2.4142 \ldots) \\
& =-0.8813 \ldots+0.8813 \ldots=0
\end{aligned}
\]
hvilket er i orden.

\[
\begin{aligned}
& x=-5.4142 \ldots: \\
& \ln (-5.4142 \ldots+3)+\ln (-5.4142 \ldots+5) \\
& =\ln (-2.4142 \ldots)+\ln (-0.4142 \ldots)
\end{aligned}
\]

hvilket ikke er defineret.

Den eneste gyldige løsning er $x=-2.5857 \ldots$.
\end{solution}
    
\end{example}



